% Part: second-order-logic
% Chapter: syntax-and-semantics
% Section: expressive-power

\documentclass[../../../include/open-logic-section]{subfiles}

\begin{document}

\olfileid{sol}{syn}{exp}
\olsection{અભિવ્યક્તિની શક્તિ}

\begin{explain}
દ્વિતીય-ક્રમ ચલો પરનું પરિમાણીકરણ ભાષાની અભિવ્યક્તિની શક્તિને
પ્રથમ-ક્રમ તર્કશાસ્ત્રની સરખામણીએ ઘણું વધારે છે. દ્વિતીય-ક્રમ
અસ્તિત્વવાચી પરિમાણીકરણથી અમુક ગુણધર્મો ધરાવતાં વિધેયો કે સંબંધો
અસ્તિત્વમાં છે એવું કહી શકાય છે. પ્રથમ-ક્રમ તર્કશાસ્ત્રમાં એવું કહેવા
માટે તે હેતુનો અતાર્કિક સંકેત, એટલે !!a{function} અથવા !!{predicate},
નિર્દિષ્ટ કરવો પડે છે. દ્વિતીય-ક્રમ સર્વવાચી પરિમાણીકરણથી વ્યાપકક્ષેત્રના
બધા ઉપગણો, તેના પરના બધા સંબંધો અથવા !!{domain}માંથી !!{domain}માં
જતાં બધાં વિધેયો અમુક ગુણધર્મ ધરાવે છે એવું કહી શકાય છે. પ્રથમ-ક્રમ
તર્કશાસ્ત્રમાં માત્ર ભાષાના કોઈ અતાર્કિક સંકેતને અપાયેલા ઉપગણો,
સંબંધો કે વિધેયો તે ગુણધર્મ ધરાવે છે એવું કહી શકાય. વળી,
અમુક ગુણધર્મવાળા ઉપગણો, સંબંધો કે વિધેયો અસ્તિત્વમાં છે અથવા તે બધા
એ ગુણધર્મ ધરાવે છે એમ કહીએ ત્યારે તે ગુણધર્મ નિર્દિષ્ટ કરવામાં પણ
દ્વિતીય-ક્રમ પરિમાણીકરણ વાપરી શકાય છે. તેથી પ્રથમ-ક્રમ તર્કશાસ્ત્રમાં
અવ્યાખ્યેય સંબંધો વ્યાખ્યાયિત કરી શકાય છે અને વ્યાપકક્ષેત્રના એવા
ગુણધર્મો વ્યક્ત કરી શકાય છે જે પ્રથમ-ક્રમ તર્કશાસ્ત્રમાં વ્યક્ત થતા
નથી.
\end{explain}

\begin{defn}
જો~$\Struct{M}$ એ ભાષા~$\Lang{L}$ માટેની !!a{structure} હોય, તો
સંબંધ~$R \subseteq \Domain{M}^2$ એ~$\Lang{L}$માં \emph{વ્યાખ્યેય}
છે જો એવું કોઈ !!{formula}~$!A_R(x, y)$ હોય જેમાં માત્ર $x$ અને~$y$
મુક્ત હોય અને $s(x) = a$ તથા $s(y) = b$ માટે $R(a, b)$, એટલે
$\tuple{a,b} \in R$, ત્યારે અને માત્ર ત્યારે જ જ્યારે
$\Sat{M}{!A_R(x, y)}[s]$.
\end{defn}

\begin{ex}
પ્રથમ-ક્રમ તર્કશાસ્ત્રમાં અભિન્નતા સંબંધ~$\Id{\Domain{M}}$, એટલે
$\Setabs{\tuple{a,a}}{a \in
  \Domain{M}}$, સૂત્ર~$\eq[x][y]$ વડે વ્યાખ્યાયિત કરી શકાય છે.
દ્વિતીય-ક્રમ તર્કશાસ્ત્રમાં આ સંબંધ~\emph{$\eq$ વિના} વ્યાખ્યાયિત
કરી શકાય છે. જો~$a$ અને~$b$ એ~$\Domain{M}$નો એક જ !!{element}
હોય, તો તેઓ~$\Domain{M}$ના એ જ ઉપગણોના !!{element}s હોય છે (કારણ
કે ગણો તેમના !!{element}s વડે નક્કી થાય છે). ઊલટું, જો~$a$ અને~$b$
જુદા હોય, તો તેઓ એકસરખા ઉપગણોના !!{element}s નથી: દાખલા તરીકે,
$a \neq b$ હોય તો $a \in \{a\}$ પરંતુ $b
\notin \{a\}$. તેથી ``$\Domain{M}$ના એ જ ઉપગણોના !!{element}s
હોવું'' એવો સંબંધ~$a$ અને~$b$ માટે ત્યારે અને માત્ર ત્યારે જ સાચો
છે જ્યારે $a = b$. $\Domain{M}$ના બધા ઉપગણો પર પરિમાણીકરણ કરી
શકાતું હોવાથી આ સંબંધ દ્વિતીય-ક્રમ તર્કશાસ્ત્રમાં વ્યક્ત કરી શકાય છે.
આમ, નીચેનું
!!{formula}~$\Id{\Domain{M}}$ને વ્યાખ્યાયિત કરે છે:
\[
\lforall[X][(X(x) \liff X(y))]
\]
\end{ex}

\begin{prob}
બતાવો કે $\lforall[X][(X(x) \lif X(y))]$ પણ~$\Id{\Domain{M}}$
વ્યાખ્યાયિત કરે છે (ધ્યાન આપો: અહીં $\lif$ છે, $\liff$ નહિ!).
\end{prob}

\begin{ex}
જો~$R$ દ્વિસ્થાની !!{predicate} હોય, તો~$\Assign{R}{M}$ એ
$\Domain{M}$ પરનો દ્વિસ્થાની સંબંધ છે. થોડી ગૂંચવણ થવાની શક્યતા
હોવા છતાં, આપણે~$R$નો ઉપયોગ !!{predicate}~$R$ના સંકેત માટે અને
સંબંધ~$\Assign{R}{M}$ પોતે દર્શાવવા માટે પણ કરીશું. $R$નું
\emph{પરંપરિત સંવરણ}~$R^*$ એવો સંબંધ છે કે~$a$ અને~$b$ વચ્ચે તે
ત્યારે અને માત્ર ત્યારે જ હોય જ્યારે અમુક $c_1$, \dots, $c_k$
માટે $R(a,c_1)$, $R(c_1, c_2)$, \dots, $R(c_k,b)$ સાચાં હોય.
તેમાં $k = 0$નો કિસ્સો પણ છે: $R(a,b)$ હોય તો $R^*(a,b)$ પણ હોય.
તેથી $R \subseteq R^*$. હકીકતમાં~$R^*$ એવો સૌથી નાનો સંબંધ છે જે
$R$ને સમાવે છે અને પરંપરિત છે. દ્વિતીય-ક્રમ તર્કશાસ્ત્રમાં~$X$ એ~$R$
સમાવતો પરંપરિત સંબંધ છે એવું કહી શકાય છે:
\begin{multline*}
  !B_R(X) \ident \lforall[x][\lforall[y][(R(x,y) \lif X(x, y))]] \land {}\\
\lforall[x][\lforall[y][\lforall[z][((X(x,y) \land X(y,z)) \lif X(x,
      z))]]].
\end{multline*}
પહેલો સંયોજક $R \subseteq X$ કહે છે અને બીજો~$X$ પરંપરિત છે એમ કહે
છે.

$X$ એવો સૌથી નાનો સંબંધ છે એમ કહેવાનો અર્થ એ છે કે~$R$ને સમાવતા
અને પરંપરિત એવા દરેક સંબંધમાં તે પોતે સમાવિષ્ટ છે. તેથી~$R$નું
પરંપરિત સંવરણ નીચેના !!{formula} વડે વ્યાખ્યાયિત કરી શકાય છે:
\[
R^*(X) \ident !B_R(X) \land \lforall[Y][(!B_R(Y) \lif
  \lforall[x][\lforall[y][(X(x, y) \lif Y(x,y))]])].
\]
$\Sat{M}{R^*(X)}[s]$ ત્યારે અને માત્ર ત્યારે જ જ્યારે $s(X) = R^*$.
$R$નું પરંપરિત સંવરણ પ્રથમ-ક્રમ તર્કશાસ્ત્રમાં વ્યક્ત કરી શકાતું નથી.
\end{ex}

\end{document}
